Oljer har gjort så mycket matte att han har tröttnat på primtal. 
Han har nu börjat studera så kallade trevliga tripletter. 
Det är väldigt lätt att skapa en trevlig triplett. 
Börja med ett heltal $a$ som är
större än noll och skapa sedan tripletten $(a,a+1,a+2)$. 
Alla tripletter som skapats på detta sätt är trevliga. 
Några exempel på trevliga tripletter är $(4,5,6)$ och $(15,16,17)$.

För att lära sig mer om trevliga tripletter vill han hitta hur 
många tal som kan skrivas som produkten av
alla talen i en trevlig triplett. Han kallar dessa "treiga tal". 
Några exempel på treiga tal är $24=2 \cdot 3 \cdot 4$
och $336 = 6 \cdot 7 \cdot 8$. Han vill nu veta hur många treiga tal det 
finns som är mindre än talet $n$.

\section*{Indata}
Den första raden innehåller heltalet $n$ ($1 \le n \le 10^9$), som beskrivits ovan.


\section*{Utdata}
Skriv ut ett heltal: antalet treiga tal mindre än $n$.


\section*{Poängsättning}
Din lösning kommer att testas på en mängd testfallsgrupper.
För att få poäng för en grupp så måste du klara alla testfall i gruppen.

\noindent
\begin{tabular}{| l | l | p{12cm} |}
  \hline
  \textbf{Grupp} & \textbf{Poäng} & \textbf{Gränser} \\ \hline
  $1$    & $20$       & $n \leq 30$ \\ \hline
  $2$    & $20$       & $n \leq 80$ \\ \hline
  $3$    & $20$       & $n \leq 1000$ \\ \hline
  $4$    & $20$       & $n \leq 10^5$ \\ \hline
  $5$    & $20$       & Inga ytterligare begränsningar. \\ \hline
\end{tabular}
